\documentclass[11pt,a4paper]{article}

% --- Essential Typographic & Layout Packages ---
\usepackage[utf8]{inputenc}
\usepackage[T1]{fontenc}
\usepackage[margin=1in]{geometry} % Professional page margins
\usepackage{microtype}            % Subliminal typographic micro-adjustments
\usepackage{amsmath,amsfonts}     % Advanced math typesetting
\usepackage{ulem}                 % Needed for strike-through (\sout)
\usepackage{enumitem}             % Customized lists without ugly gaps

% --- Color & UI Background Packages ---
\usepackage[dvipsnames]{xcolor}
\usepackage[most]{tcolorbox}      % For premium callout boxes / backgrounds

% --- Design Palette Definition ---
\definecolor{DeepNavy}{RGB}{24, 43, 73}
\definecolor{AccentBlue}{RGB}{54, 95, 145}
\definecolor{LightBg}{RGB}{245, 247, 250}
\definecolor{BorderGray}{RGB}{218, 224, 233}

% --- Custom Header Formatting ---
\newcommand{\Heading}[1]{\vspace{0.4cm}\noindent\textbf{\textcolor{DeepNavy}{\large #1}}\par\vspace{0.15cm}}

\begin{document}

% --- Title Section ---
\begin{center}
    {\LARGE\bfseries\textcolor{DeepNavy}{Advanced Calculus}}
    \vspace{0.6cm}
\end{center}

\Heading{Differential Equations (DE):}
An equation involving derivatives of one or more dependent variables with respect to one or more independent variables is called a differential equation.

\Heading{Ordinary Differential Equations (ODE):}
A differential equation involving ordinary derivatives of one or more dependent variables with respect to a single independent variable is called an ordinary differential equation.

\Heading{Partial Differential Equation (PDE):}
A differential equation involving partial derivatives of one or more dependent variables with respect to more than one independent variable is called a partial differential equation.

% --- Stylized Background Callout for Examples ---
\begin{tcolorbox}[colback=LightBg, colframe=BorderGray, arc=4mm]
    {\small\textbf{\textcolor{AccentBlue}{EXAMPLES BY CATEGORY}}}
    \vspace{0.2cm}
    \begin{description}[leftmargin=0.4cm, style=nextline, font=\normalfont\bfseries\textcolor{DeepNavy}]
        
        \item[Ordinary Differential Equations (ODE)]
        \[ \frac{d^2y}{dx^2} + xy\left(\frac{dy}{dx}\right)^2 = 0 \]
        \[ \frac{d^4x}{dt^4} + 5\frac{d^2x}{dt^2} + 3x = \sin(t) \]
        
        \item[Partial Differential Equations (PDE)]
        \[ \frac{\partial v}{\partial \sigma} + \frac{\partial v}{\partial t} = v \]
        \[ \frac{\partial^2u}{\partial x^2} + \frac{\partial^2u}{\partial y^2} + \frac{\partial^2u}{\partial z^2} = 0 \]
        
        \item[Not DE but Basic Derivatives]
        \[ \frac{d}{dx}\left(e^{ax}\right) = ae^{ax} \]
        \[ \frac{d}{dx}(uv) = u\frac{dv}{dx} + v\frac{du}{dx} \]
        
    \end{description}
\end{tcolorbox}

\Heading{Order of a DE:}
The order of the highest-order derivative involved in a differential equation is called the order of the differential equation.

\Heading{Degree of a DE:}
The degree of a differential equation is the power of the highest-order derivative present in the equation.

\begin{tcolorbox}[colback=LightBg, colframe=BorderGray, arc=4mm]
    {\small\textbf{\textcolor{AccentBlue}{EXAMPLES BY CATEGORY}}}
    \vspace{0.2cm}
    \begin{description}[leftmargin=0.4cm, style=nextline, font=\normalfont\bfseries\textcolor{DeepNavy}]
        \item[1st Order 1st Degree ODE]
        \[ \frac{dy}{dx} + y = e^x \]
        
        \item[2nd Order 1st Degree ODE]
        \[ \frac{d^2y}{dx^2} + 3\frac{dy}{dx} + 2y = 0 \]
        
        \item[2nd Order 2nd Degree ODE]
        \[ \left(\frac{d^2y}{dx^2}\right)^2 + y\frac{dy}{dx} = 0 \]
        
        \item[1st Order 1st Degree PDE]
        \[ \frac{\partial u}{\partial t} + \frac{\partial u}{\partial x} = 0 \]
        
        \item[2nd Order 1st Degree PDE]
        \[ \frac{\partial^2 u}{\partial t^2} - c^2 \frac{\partial^2 u}{\partial x^2} = 0 \]
        
        \item[2nd Order 2nd Degree PDE]
        \[ \left(\frac{\partial^2 u}{\partial x^2}\right)^2 + \left(\frac{\partial^2 u}{\partial y^2}\right)^2 = 0 \]
    \end{description}
\end{tcolorbox}

\Heading{Linear ODE:}
A linear ordinary differential equation of order $n$, in the dependent variable $y$ and the independent variable $x$, is an equation that can be expressed in the form
\[
a_0 y^{(n)} + a_1 y^{(n-1)} + \ldots + a_n y = f(x)
\]
or equivalently,
\[
a_0(x) \frac{d^n y}{dx^n} + a_1(x) \frac{d^{n-1} y}{dx^{n-1}} + \ldots + a_n(x) y = f(x)
\]
where $a_0, a_1, \ldots, a_n$ are functions of $x$ (not identically zero) and $f(x)$ is a given function.

\vspace{0.2cm}
\noindent\textbf{\textcolor{DeepNavy}{Important Observations:}}
\begin{enumerate}[label=\arabic*), leftmargin=0.6cm, itemsep=0.05cm]
    \item The dependent variable $y$ and its various derivatives occur to the first degree only.
    \item No products of $y$ and/or any of its derivatives are present.
    \item No transcendental functions of $y$ or its derivatives occur (e.g., $\sin y$, $e^y$).
\end{enumerate}

\Heading{Nonlinear ODE:}
A nonlinear ordinary differential equation is an ordinary differential equation that is not linear.

\begin{tcolorbox}[colback=LightBg, colframe=BorderGray, arc=4mm]
    {\small\textbf{\textcolor{AccentBlue}{EXAMPLES BY CATEGORY}}}
    \vspace{0.2cm}
    \begin{description}[leftmargin=0.4cm, style=nextline, font=\normalfont\bfseries\textcolor{DeepNavy}]
        \item[Linear ODE with Constant Coefficients]
        \[ \frac{d^2y}{dx^2} + 3\frac{dy}{dx} + 2y = 0 \]
        
        \item[Linear ODE with Variable Coefficients]
        \[ \frac{d^4y}{dx^4} + x^{30}\frac{dy}{dx} + 2y = xe^x \]
        
        \item[Nonlinear ODE Examples]
        \[ \left(\frac{d^2y}{dx^2}\right)^2 + y\frac{dy}{dx} = 0 \]
        \[ \frac{dy}{dx} + \sin(y) = 0 \] % Replaced the incorrect linear equation with a proper nonlinear one
        \[ A \frac{d^2y}{dx^2} + K \frac{d^3y}{dx^3} + x \frac{dy}{dx} = x e^x \frac{dx}{dy} \]
        
        \item[General Form of a Nonlinear ODE]
        \[ F\left(x, y, \frac{dy}{dx}, \frac{d^2y}{dx^2}, \ldots, \frac{d^ny}{dx^n}\right) = 0 \]
    \end{description}
\end{tcolorbox}

\Heading{What Makes an ODE Nonlinear?}
An ODE is classified as nonlinear if it violates any of the key linear criteria:
\begin{itemize}[leftmargin=0.5cm, itemsep=0.05cm]
    \item \textbf{Higher Powers:} Any $y$ or derivative term has an exponent other than 1.
    \item \textbf{Product Terms:} The variable $y$ is multiplied by its own derivatives.
    \item \textbf{Nonlinear Functions:} The equation contains functions like $\sin(y)$, $e^{y}$, or $\ln(y)$.
\end{itemize}

\Heading{Origin and Application:}
Differential equations occur in connection with numerous problems encountered across various branches of science and engineering:
\begin{itemize}[leftmargin=0.5cm, itemsep=0.05cm]
    \item Population Dynamics
    \item Dynamical Systems and Chaos
    \item Fluid Dynamics
    \item Heterogeneous Heat Transfer
    \item Biophysics and Mathematical Biology
\end{itemize}

\Heading{Formation of Differential Equations:}
To form a differential equation from a given relation involving independent/dependent variables and arbitrary constants, we eliminate the arbitrary constants by successive differentiation.

\noindent\textbf{Example Problem:} \\
Form the differential equation by eliminating the constants $a$ and $b$ from:
\begin{equation}
    y = ax + bx^2
\end{equation}

\noindent\textbf{Solution:} \\
Given the relation with two arbitrary constants $a$ and $b$:
\begin{equation}
    y = ax + bx^2 \label{eq:original}
\end{equation}

Since there are two arbitrary constants, we differentiate the equation twice with respect to $x$ to eliminate them.

\paragraph{Step 1: First Differentiation}
Differentiating Equation \eqref{eq:original} with respect to $x$:
\begin{equation}
    \frac{dy}{dx} = a + 2bx \label{eq:diff1}
\end{equation}

\paragraph{Step 2: Second Differentiation}
Differentiating Equation \eqref{eq:diff1} with respect to $x$:
\begin{equation}
    \frac{d^2y}{dx^2} = 2b \implies b = \frac{1}{2}\frac{d^2y}{dx^2} \label{eq:b_val}
\end{equation}

\paragraph{Step 3: Eliminate Constant $a$}
Substitute the value of $b$ from Equation \eqref{eq:b_val} into Equation \eqref{eq:diff1}:
\[
\frac{dy}{dx} = a + 2\left(\frac{1}{2}\frac{d^2y}{dx^2}\right)x
\]
\[
\frac{dy}{dx} = a + x\frac{d^2y}{dx^2} \implies a = \frac{dy}{dx} - x\frac{d^2y}{dx^2} \label{eq:a_val}
\]

\paragraph{Step 4: Substitute $a$ and $b$ into the Original Equation}
Substitute the expressions for $a$ and $b$ back into Equation \eqref{eq:original}:
\[
y = \left(\frac{dy}{dx} - x\frac{d^2y}{dx^2}\right)x + \left(\frac{1}{2}\frac{d^2y}{dx^2}\right)x^2
\]
Expand the terms:
\[
y = x\frac{dy}{dx} - x^2\frac{d^2y}{dx^2} + \frac{1}{2}x^2\frac{d^2y}{dx^2}
\]
Combine the like terms of $\frac{d^2y}{dx^2}$:
\[
y = x\frac{dy}{dx} - \frac{1}{2}x^2\frac{d^2y}{dx^2}
\]

\paragraph{Final Form:}
Multiply the entire equation by $2$ and rearrange terms to obtain the standard form:
\begin{equation}
    x^2\frac{d^2y}{dx^2} - 2x\frac{dy}{dx} + 2y = 0
\end{equation}
This is the required ordinary differential equation of \text{$2^{nd}$ order and $1^{st}$ degree.}


\end{document}
